% =============================================================
% Annexure K: Glossary of Key Terms
% =============================================================
\section{Glossary of Key Terms}
\label{annexure:glossary}

The following definitions establish the meaning of key terms as they are
used throughout this dissertation.  Formal mathematical definitions are
provided in the relevant chapters.

\begin{description}
    \item[Deep neural network (DNN).]  A feed-forward artificial neural
          network with multiple hidden layers, each consisting of an
          affine transformation followed by a nonlinear activation
          function.  In this dissertation, the DNN is used as a
          regression model to predict continuous insurance costs from
          policyholder attributes.

    \item[Funnel architecture.]  A network architecture in which the
          number of neurons decreases monotonically from the input layer
          to the output layer.  The specific architecture used in this
          study has layer widths $256 \to 128 \to 64 \to 32 \to 16$.

    \item[Generalised linear model (GLM).]  A class of statistical
          models that generalise ordinary linear regression by allowing
          the response variable to follow any distribution in the
          exponential family and by relating the mean of the response to
          the linear predictor through a link function.  In this
          dissertation, a Gamma GLM with a logarithmic link function
          serves as the actuarial benchmark.

    \item[Random forest.]  An ensemble learning method that constructs a
          collection of decision trees, each trained on a bootstrap
          sample of the data with a random subset of features considered
          at each split.  The final prediction is the average of the
          individual tree predictions.  In this dissertation, a random
          forest regressor serves as one of three benchmark models
          against which the DNN is compared.

    \item[Gradient-boosted trees (XGBoost).]  An ensemble method that
          builds an additive sequence of shallow decision trees, with
          each new tree fitted to the residuals of the preceding
          ensemble.  XGBoost \citep{chen2016xgboost} is a widely used
          implementation that incorporates regularisation, approximate
          histogram-based splitting and early stopping.  In this
          dissertation, an XGBoost regressor serves as one of three
          benchmark models.

    \item[Monte Carlo simulation.]  A computational method that uses
          repeated random sampling from a specified probability
          distribution to estimate the properties of a quantity of
          interest.  In this study, Monte Carlo simulation generates
          synthetic policyholder profiles from fitted parametric
          marginals and propagates them through the trained network to
          obtain a predicted-cost distribution.

    \item[Model-implied predicted-cost distribution.]  The distribution
          of predicted costs obtained by evaluating the trained
          prediction function at a large sample of input profiles drawn
          from a specified input distribution.  This distribution
          characterises the range, central tendency, dispersion and shape
          of the model's output under population-level input variation.

    \item[Shapley value.]  A concept from cooperative game theory that
          assigns to each player in a coalition game a unique payoff
          satisfying the axioms of efficiency, symmetry, linearity and
          the null-player property.  In the context of model
          interpretation, the ``players'' are input features and the
          ``payoff'' is the prediction for a given observation.  The
          Shapley value quantifies the average marginal contribution of
          each feature across all possible feature coalitions.

    \item[SHAP (SHapley Additive exPlanations).]  A framework that
          operationalises the computation of approximate Shapley values
          for machine-learning models.  The gradient-based variant used
          in this dissertation exploits the differentiability of the
          neural network to compute Shapley-value approximations
          efficiently.

    \item[Conditional Shapley analysis.]  The computation and
          summarisation of Shapley values within a subset of predictions
          that exceed a specified threshold.  In this dissertation,
          conditional Shapley summaries are computed at the 90th, 95th
          and 99th percentiles of the predicted-cost distribution.

    \item[Feature-importance ranking reversal.]  The phenomenon whereby
          the ordering of features by their mean absolute Shapley values
          differs between the global summary and the conditional summary
          at a specific threshold.  A ranking reversal indicates that the
          variables driving predictions in the specified region differ
          from those that dominate on average.

    \item[Data-leakage prevention.]  The set of protocols that ensure no
          information from the validation or test partitions contaminates
          the training process.  In this dissertation, all preprocessing
          parameters are computed exclusively on the training partition
          and applied to the validation and test partitions without
          refitting.

    \item[Residual-augmented sensitivity analysis.]  A diagnostic
          procedure that adds scaled fitted residuals to the Monte Carlo
          predicted costs to examine whether the distributional
          conclusions are robust to the residual structure of the fitted
          model.
\end{description}
